%-----------------------------------------------------------
% Model Overview
%-----------------------------------------------------------

The mannequin geometry was processed to create a shell with a uniform wall thickness of 3.2 mm in the area encompassing the 21 probes according to the the ten-twenty electrode system of the International Federation~\cite{klem-1999}. 
This shell was divided into two interlocking halves (superior and inferior) to simplify printing, 
assembly, and internal inspection.

The internal cavity was designed to be filled with a gelatine-salt-water solution. This space establishes the physical domain where the internal antennas and wire guides are embedded, allowing the propagation of electrical fields to be studied under controlled conditions (see Section~\ref{sec:gelatine-props}). \nota{la seccion ref aca deberia ser una seccion sobre la gelatina, sus propiedades, composicion, etc pero todavia no existe, chequear despues}

Dedicated openings were integrated into the lid to facilitate gel pouring and the alignment 
of the internal wire-guides.

\begin{figure}[h]
    \centering
    \includegraphics[width=0.65\textwidth]{figures/head_shell_overview.png}
    \caption{Overall 3D model of the phantom head assembly, highlighting major components.}
    \label{fig:head_shell_overview}
\end{figure}

\nota{INCLUIR IMAGEN DE TODO EL MODELO DESPUES. Tengo que hacerle una base para que no se caiga y tenga lugar abajo de la lid y arriba de la mesa para sacar los cables por adentro.}
